% ===================================================================
\section{Qué es la aplicación}

La aplicación de PULSO es un laboratorio interactivo. Permite construir una red de
pagos, aplicarle un choque y ver cómo se propaga, sin escribir código. También muestra
los resultados del proyecto sobre riesgo de crédito.

Tiene diez páginas en dos bloques:

\begin{center}
\small
\begin{tabular}{p{0.32\textwidth}p{0.6\textwidth}}
\toprule
\textbf{Páginas} & \textbf{Para qué sirven} \\
\midrule
Inicio & Presentación y carga del escenario de ejemplo. \\
1 Construir red & Crear la red de pagos con la que trabajan las páginas 2 a 7. \\
2 Diagnóstico & Saber si la red amortigua o amplifica un choque. \\
3 Ritmo de pagos & Ver cómo se sincronizan los ciclos de pago. \\
4 Simulador de choques & Aplicar un choque y ver la cascada de impagos. \\
5 Explorar escenarios & Consultar 20\,000 simulaciones ya hechas. \\
6 Mapa de riesgo & Predicción instantánea del riesgo de cada entidad. \\
7 Acerca del modelo & Ecuaciones, supuestos y límites. \\
\midrule
8 Cartera multicapa & Clientes de un banco en tres capas de pagos (tarea 6). \\
9 Riesgo crediticio híbrido & Comparar el score clásico con el modelo que usa la red
(tarea 6). \\
\bottomrule
\end{tabular}
\end{center}

Las páginas 1 a 7 trabajan con una \textbf{red de entidades}: bancos, empresas o
participantes de un sistema de pagos, sin distinguir su tipo. Las páginas 8 y 9 trabajan
con una \textbf{cartera de clientes} de un banco: personas y empresas.

\begin{aviso}
Todos los datos son simulados y los parámetros son ilustrativos. La aplicación sirve
para entender mecanismos, no para medir el riesgo de una institución real.
\end{aviso}

% ===================================================================
\section{Cómo abrir la aplicación}

\subsection{En el computador, con uv}
En una terminal, dentro de la carpeta del proyecto:
\begin{verbatim}
uv sync                              # solo la primera vez: instala todo
uv run streamlit run app/Inicio.py   # abre la aplicación
\end{verbatim}
El navegador se abre en \texttt{http://localhost:8501}. Si ese puerto está ocupado,
use otro puerto y abra \texttt{http://localhost:8502}:
\begin{verbatim}
uv run streamlit run app/Inicio.py --server.port=8502
\end{verbatim}

\subsection{Con Docker}
\begin{verbatim}
docker build -t pulso .
docker run -p 8501:8501 pulso
\end{verbatim}
y abra \texttt{http://localhost:8501}. La imagen ya trae los datos y los modelos
entrenados.

\subsection{Para cerrarla}
Cierre la pestaña del navegador y, en la terminal, presione \texttt{Ctrl+C}.

% ===================================================================
\section{Elementos de la pantalla}

\begin{description}
\item[Barra lateral.] A la izquierda. Tiene la lista de páginas y, debajo, un resumen de
la red activa (topología, número de entidades y aristas) o de la cartera activa
(personas, firmas, informalidad y semilla). Se oculta con la flecha \texttt{<<}.
\item[Deslizador.] Barra con un punto que se arrastra para elegir un número. El valor
actual aparece encima. Con el teclado, las flechas lo mueven paso a paso.
\item[Formulario.] Recuadro con varios controles y un botón azul (\boton{Construir},
\boton{Simular}, \boton{Simular cartera}). Los cambios dentro del recuadro no tienen
efecto hasta presionar el botón. Así se pueden cambiar varios valores a la vez sin
esperar una simulación por cada uno.
\item[Control fuera de formulario.] Se aplica al soltarlo: la página se recalcula de
inmediato.
\item[Métrica.] Número grande con su nombre encima. Algunas tienen un símbolo
\textbf{?}: al pasar el cursor muestra una explicación. Debajo puede aparecer una
flecha con la diferencia respecto a un caso de referencia.
\item[Gráfica interactiva.] Al pasar el cursor sobre un punto se ve su información. Con
el ratón se puede hacer zoom arrastrando un rectángulo; doble clic restablece la vista.
En la esquina superior derecha hay botones para descargar la imagen, mover, ampliar y
ver en pantalla completa. Al hacer clic en un elemento de la leyenda se oculta o se
muestra esa serie.
\item[Tabla.] Al hacer clic en el nombre de una columna se ordena por ella.
\item[Desplegable (\emph{expander}).] Línea con una flecha \texttt{>} que se abre al
hacer clic y muestra contenido adicional.
\item[Indicador de ejecución.] Mientras la aplicación calcula, en la esquina superior
derecha aparece un aviso de ejecución. Espere a que desaparezca antes de leer los
resultados.
\end{description}

\subsection{Estado compartido entre páginas}
La red que se construye en \emph{Construir red} y los parámetros que se ajustan en
\emph{Diagnóstico}, \emph{Ritmo de pagos} y \emph{Simulador de choques} se comparten.
Por ejemplo, el valor de $\alpha$ elegido en \emph{Diagnóstico} se usa también en el
simulador. Lo mismo ocurre con la cartera de la página 8, que usa la página 9.

El estado dura mientras la pestaña esté abierta. Si se recarga la página del navegador,
todo vuelve a los valores iniciales.

\subsection{Límite de tamaño}
Con más de 300 entidades, las páginas que simulan en vivo (\emph{Diagnóstico}, en su
choque de prueba; \emph{Ritmo de pagos}; \emph{Simulador de choques}; y la verificación
del \emph{Mapa de riesgo}) se desactivan y muestran un aviso. La barra lateral también lo
indica. El \emph{Mapa de riesgo} sigue funcionando porque predice sin simular.

% ===================================================================
\section{Conceptos mínimos}

Estos términos aparecen en todas las páginas.

\begin{description}
\item[Entidad o nodo.] Un participante de la red. Se dibuja como un círculo.
\item[Arista.] Una flecha de pago de una entidad a otra. Su peso es el pago medio por
ciclo.
\item[Ciclo.] Unidad de tiempo de las páginas 1 a 7: un ciclo de pagos, equivalente a un
día hábil. En las páginas 8 y 9 la unidad es el mes.
\item[Pago medio.] Unidad de dinero de las páginas 1 a 7. Los montos están escalados para
que la arista promedio valga 1.
\item[Déficit ($x$).] Faltante de liquidez de una entidad. Cero significa sin problemas.
\item[Colchón (umbral).] Déficit máximo que una entidad soporta. Si lo supera, entra en
impago.
\item[Impago.] La entidad deja de pagar a todos sus receptores.
\item[Cascada.] Impagos en cadena: el impago de una entidad lleva a otras a superar su
colchón.
\item[Fase ($\theta$).] Momento del ciclo de pagos en que está una entidad, como la
posición de la manecilla de un reloj.
\item[Sincronía ($r$).] Qué tan alineados están los ciclos. $r=1$: todos pagan al mismo
tiempo. $r=0$: cada uno a su ritmo.
\item[Semilla.] Número que fija el azar. Con la misma semilla se obtiene exactamente el
mismo resultado.
\item[PD.] Probabilidad de impago de un cliente en los próximos doce meses.
\end{description}
